% Generated draft: requires mathematical and source review.
\documentclass[11pt]{article}
\usepackage[margin=25mm]{geometry}
\usepackage{fontspec}
\setmainfont{Latin Modern Roman}
\usepackage{amsmath,amssymb,mathtools,unicode-math}
\setmathfont{Latin Modern Math}
\usepackage{enumitem,needspace,xurl,hyperref}
\setlength{\parindent}{0pt}
\setlength{\parskip}{5pt}
\setlength{\emergencystretch}{3em}
\newcommand{\N}{\mathbb{N}}
\newcommand{\Z}{\mathbb{Z}}
\newcommand{\Q}{\mathbb{Q}}
\newcommand{\R}{\mathbb{R}}
\newcommand{\C}{\mathbb{C}}
\newcommand{\F}{\mathbb{F}}
\newcommand{\A}{\mathbb{A}}
\newcommand{\PP}{\mathbb{P}}
\newcommand{\E}{\mathbb{E}}
\newcommand{\eps}{\varepsilon}
\newcommand{\dd}{\,\mathrm{d}}
\begin{document}
\section*{UnsolvedMath ID: 20003351}
\textbf{Source problem code:} catalogue entry
\subsection{Definitions and mathematical statement}
The exact source item is a historical continuation of Conjecture 51: it says that the one-dimensional-Zariski-closure case was proved assuming ``generalized Riemannian hypothesis'', citing [203]. The supplied aid identifies that phrase as a typo for the Generalized Riemann Hypothesis (GRH) and identifies the subject as Andr\'e--Oort. The original conjecture and its ambient Shimura variety are not present in this fragment.

For context, the Andr\'e--Oort assertion is that in a Shimura variety the irreducible components of the Zariski closure of a set of special points are special subvarieties, with ``special'' interpreted in its Shimura-theoretic sense. Thus a one-dimensional closure can be reducible: its curve components are treated by an irreducible-curve theorem, while any zero-dimensional components are special points from the original set. The aid reports [203] as Yafaev's conditional irreducible-curve theorem and separately reports the full general theorem in Pila--Shankar--Tsimerman. These historical and later claims are distinct from the supplied attempt's own contribution and require source verification.

The aid's explicit-family statement concerns only $Y(1)^n$, with the modular $j$-invariant $j:\mathbb H\to\C$, where $\mathbb H=\{\tau\in\C:\operatorname{Im}\tau>0\}$. Fix
\[
 n\geq2,\qquad N_1=1,\qquad N_2,\ldots,N_n\in\Z_{\geq2}.
\]
Let $\Phi_N(X,Y)$ be the classical modular polynomial for cyclic $N$-isogenies and define
\[
 V=\{(x_1,\ldots,x_n)\in\C^n:
             \Phi_{N_r}(x_1,x_r)=0\ (2\leq r\leq n)\},
\]
\[
 C=\overline{\{(j(\tau),j(N_2\tau),\ldots,j(N_n\tau)):
                 \tau\in\mathbb H\}}^{\mathrm{Zar}}.
\]
The claimed developed result is that $C$ is an irreducible special curve and an irreducible component of $V$, and that
\[
 \{(j(f\mathrm i),j(N_2f\mathrm i),\ldots,j(N_nf\mathrm i)):
                        f\in\Z_{\geq1}\}
\]
consists of pairwise distinct special points Zariski dense in $C$. Here $\mathrm i^2=-1$ is fixed, so these are CM points from the fixed imaginary quadratic field $\Q(\mathrm i)$ with increasing conductors; $f$ is an integer parameter, not an index on $j$. The common parameter $\tau$ selects a particular component of the simultaneous modular-polynomial intersection. The statement does not identify all of $V$ as a single irreducible curve or claim independent coordinates.

The registry labels the record solved. The supplied aid describes the full Andr\'e--Oort problem as solved in the cited later literature, while classifying its own explicit-family work as a partial result. Neither label is a claim of independent mathematical review here.
\subsection{Short English statement}
The fragment records a historical GRH-dependent curve case of Andr\'e--Oort; a supplied later reference is reported to prove the general theorem. Its separate explicit example is a fixed-field CM conductor ray, Zariski dense in the special curve selected by a common modular parameter, rather than in every component of the modular-polynomial intersection.
\subsection{Sources}
\begin{itemize}
\item Local UnsolvedMath registry, permanent ID 20003351, exact source code catalogue entry; \texttt{aim-other-notes.json}, notes[157], historical continuation of Conjecture 51.
\item American Institute of Mathematics, \emph{Emerging applications of measure rigidity}, \url{https://aimath.org/WWN/measrigid/measrigid.pdf}; numerical citation [203]. The source phrase ``generalized Riemannian hypothesis'' is preserved as a quoted typo, with the aid's GRH interpretation explicit.
\item Supplied identification of [203]: Andrei Yafaev, arXiv:math/0302125, later \emph{Duke Mathematical Journal} 132 (2006), 393--407. Supplied later general-theorem reference: Pila--Shankar--Tsimerman, arXiv:2109.08788v4. These identifications and scope claims have not been independently checked.
\item Supplied research attempt 1: partial\_result, problem status solved, rigor label proved, novelty confidence low; explicit special-curve family and reducible-component reduction. Machine-generated research aid, not peer reviewed; independently verify before citation. Original remarks and literature fields are empty.
\end{itemize}
\end{document}
